\section{Related Work}
\label{sec:related}
\paragraph{Causal skill attribution and agent evaluation.}
Paired trace auditing, no-skill comparisons, randomized skill masking, and continual skill evaluation already show that exposing an agent to a skill can change behavior and that effects are heterogeneous~\citep{zhou2026counterfactual,dong2026agentskills,wang2026notallskills,guan2026continualskillbench}. We therefore do not claim paired skill causality or heterogeneous effects as new. Our question is narrower and evaluative: does the targeted intervention improve the original agent, does the target operation output matter under the same exposure surface, and does moving the same precomputed output between two integration surfaces change the final answer under a forced one-answer harness? The observed T$>$G but T$=$N cutoff cell demonstrates why these are different estimands rather than interchangeable baselines.

\paragraph{Skill evolution and reusable procedures.}
Evo-Harness, SkillProx, HyperSkill, Anything2Skill, and related systems study continual maintenance, compilation, optimization, structured skill memory, and conversion of external knowledge into reusable procedures~\citep{wei2026evoharness,zheng2026skillprox,xu2026hyperskill,pan2026anything2skill}. These works make procedure reuse and skill accumulation established objects. We treat acquisition as upstream: after a reusable procedure exists, our contribution is an audit of the scientific credit assigned to it. A stronger control may preserve a net repair while exposing a surface interaction, eliminating a skill-container interpretation, or downgrading a previously positive support cell.

\paragraph{Temporal retrieval and context construction.}
Selective/corrective RAG and temporal retrieval already address irrelevant evidence, time-sensitive indexing, and explicit temporal constraints upstream of generation~\citep{asai2024selfrag,yan2024crag,zhang2025mrag,abdallah2026tempretriever,han2025temporalgraphs,wang2026iarag,liu2026chronos}. Consequently, a useful cutoff or alignment operation need not be packaged as callable state. Our R$_{\mathrm{surf}}$ experiment does \emph{not} compare T with an independently learned temporal retriever. It intentionally runs the same frozen operation and materializes the identical output in context, creating a same-information surface-isolation test. The result therefore speaks only to integration-surface placement under the frozen one-answer harness, not to algorithmic equivalence between skills and temporal RAG or to adaptive callable use in deployment. A broader retrieval claim would require a separately implemented retriever or reranker that must discover the relevant transformation from the same raw candidates; we do not make that claim here.

The residual novelty is thus methodological: a reusable-skill evaluation should distinguish net repair, same-surface output contrast, placebo/surface perturbation, and integration-surface value before assigning credit to ``the skill.'' In our data, making those distinctions changes multiple verdicts rather than merely adding confidence intervals to an unchanged story.
